\documentclass[11pt]{article}

\usepackage[margin=1in]{geometry}
\usepackage{amsmath,amssymb}
\usepackage{algorithm}
\usepackage{algpseudocode}
\usepackage{booktabs}
\usepackage{tabularx}
\usepackage{array}
\usepackage{xcolor}
\usepackage{microtype}
\usepackage[htt]{hyphenat}
\usepackage{hyperref}
\usepackage{enumitem}
\usepackage{parskip}
\usepackage{titlesec}
\usepackage{graphicx}

\sloppy
\emergencystretch=3em
\hbadness=10000
\hfuzz=2pt

\hypersetup{
    colorlinks=true,
    linkcolor=blue!60!black,
    citecolor=blue!60!black,
    urlcolor=blue!60!black,
}

\newcommand{\strategy}[1]{\texttt{#1}}

\title{\textbf{Masking Strategies for Antibody MLM} \\[4pt]
\large A Plain-Language Guide to Every Algorithm}
\author{}
\date{}

\begin{document}
\maketitle

\begin{abstract}
This document explains every masking strategy in our antibody masked
language model (MLM) framework. It covers \emph{what} each strategy
does, \emph{why} we designed it that way, and \emph{what alternatives
we considered}. Written to be readable without deep background in
NLP or immunology.
\end{abstract}

\tableofcontents
\newpage

% ============================================================================
\section{Background: What Is Masking and Why Does It Matter?}
% ============================================================================

In masked language modeling (MLM), we hide some tokens in a sequence
and train a model to predict them from context. The choice of
\emph{which} tokens to hide is the masking strategy. This choice
matters a lot for antibodies because different parts of an antibody
sequence carry very different amounts of information.

\subsection{The Basic Setup (Shared by All Strategies)}

Every strategy in our framework works in two stages:

\begin{enumerate}
\item \textbf{Pick positions} --- the strategy decides which residues
      to mask. This is the part that differs between strategies.
\item \textbf{Corrupt those positions} --- shared across all strategies,
      using the standard BERT recipe:
      \begin{itemize}
          \item 80\% of picked positions become \texttt{[MASK]}
          \item 10\% get replaced with a random amino acid
          \item 10\% stay unchanged (but the model still has to
                ``predict'' them)
      \end{itemize}
\end{enumerate}

This two-stage split means every strategy only has to answer one
question: \emph{which positions?} The corruption step is identical
everywhere, so comparing strategies is apples-to-apples.

\subsection{Key Parameters (Same Across All Strategies)}

\begin{itemize}
\item \textbf{Mask rate} $= 15\%$: we hide 15\% of the non-special
      tokens in every sequence, every time. Kept fixed so strategy
      comparisons are fair.
\item \textbf{Special tokens} (\texttt{[CLS]}, \texttt{[SEP]},
      \texttt{[PAD]}) are \emph{never} masked.
\item Random replacements only use the 20 canonical amino acids
      (\texttt{ACDEFGHIKLMNPQRSTVWY}), never weird tokens like
      \texttt{X} or \texttt{B}.
\end{itemize}

\subsection{The Weighted Coin-Flip Formula}
\label{sec:weighted-flip}

Most of our strategies use the same trick: give every position a
``weight'' that says how important it is to mask, then flip a
weighted coin at each position. The formula is:

\begin{equation}
\boxed{\;
\text{probability of masking position } i
\;=\;
\min\!\left(\;
0.15 \times \frac{w_i}{\text{average weight}},
\;\; 1
\;\right)
\;}
\label{eq:weighted-coin}
\end{equation}

The key insight: dividing by the average weight means that no matter
what weights you pick, you still mask about 15\% of positions overall.
High-weight positions just ``steal'' budget from low-weight ones.

\textbf{Why this formula?}
\begin{itemize}
\item It preserves the total budget: the expected number of masked
      positions is always $\approx 15\%$.
\item It separates \emph{where to focus} (the weights) from
      \emph{how much to mask} (the 15\% rate).
\item When all weights are equal, it reduces exactly to uniform
      random masking (the baseline).
\end{itemize}

% ============================================================================
\section{Strategy 1: Uniform (The Baseline)}
\label{sec:uniform}
% ============================================================================

\subsection{What It Does}
Every non-special position has an equal 15\% chance of being masked.
Nothing fancy --- this is standard BERT.

\subsection{Algorithm}
\begin{enumerate}
\item Set the masking probability to $0.15$ at every position.
\item Set it to $0$ at special tokens.
\item Flip a coin independently at each position.
\end{enumerate}

\subsection{Why We Need It}
It is our \textbf{control}. Without it, we cannot tell whether a
fancier strategy actually helps or just changes things randomly. Every
other strategy is compared against uniform to measure improvement.

% ============================================================================
\section{Strategy 2: CDR-Focused Masking}
\label{sec:cdr}
% ============================================================================

\subsection{The Problem}
An antibody has two kinds of regions:
\begin{itemize}
\item \textbf{Framework} ($\sim$75\% of residues): structurally
      conserved scaffold. Easy to predict; boring.
\item \textbf{CDRs} ($\sim$25\% of residues): hypervariable loops
      that determine what the antibody binds. Hard to predict;
      interesting.
\end{itemize}

Uniform masking spends 75\% of its budget on the easy framework
positions. That is wasteful.

\subsection{The Solution}
Give CDR positions higher masking weights so they get masked more
often. We use four weights:

\begin{center}
\begin{tabular}{lcc}
\toprule
Region & Weight & Rationale \\
\midrule
Framework & 1 & baseline \\
CDR1 & 3 & moderately variable \\
CDR2 & 3 & moderately variable \\
CDR3 & 6 & most variable, carries binding specificity \\
\bottomrule
\end{tabular}
\end{center}

These weights are plugged into the weighted coin-flip formula
(Section~\ref{sec:weighted-flip}).

\subsection{What This Achieves}
With typical antibody proportions, this puts about
\textbf{60--70\% of masks in CDRs} despite CDRs being only
$\sim$25\% of the sequence. CDR3 alone absorbs about 40\% of
the mask budget.

\textbf{Worked example.} Suppose a 120-residue VH has 90 framework
residues (weight 1), 6 CDR1 (weight 3), 8 CDR2 (weight 3), and
16 CDR3 (weight 6). The average weight is:
\[
\bar{w} = \frac{90 \times 1 + 6 \times 3 + 8 \times 3 + 16 \times 6}{120}
= \frac{90 + 18 + 24 + 96}{120} = 1.9
\]
So CDR3's masking probability is $0.15 \times 6 / 1.9 \approx 47\%$,
while framework's is $0.15 \times 1 / 1.9 \approx 8\%$.

\subsection{Algorithm}
\begin{algorithm}[H]
\caption{CDR Masking}
\begin{algorithmic}[1]
\Require CDR region labels $c_i \in \{0, 1, 2, 3\}$ for each position
\If{no CDR labels available}
    \State fall back to uniform masking
\EndIf
\State Look up weight $w_i$ from the table above using $c_i$
\State Zero out weights at special tokens
\State Compute average weight $\bar{w}$ over non-special positions
\State Set probability: $p_i = \min(0.15 \times w_i / \bar{w},\; 1)$
\State Flip coins: mask position $i$ with probability $p_i$
\end{algorithmic}
\end{algorithm}

\subsection{Design Decisions}

\textbf{Why is CDR3's weight twice CDR1/CDR2?}
CDR3 is the only CDR with V(D)J junctional diversity --- it is
far more variable and carries the bulk of binding specificity.
CDR1 and CDR2 come from the germline V gene and are less diverse.

\textbf{Why not set framework weight to 0?}
Framework residues are important for structural stability and
VH--VL pairing. Setting them to 0 would mean the model never
practices predicting them, which hurts downstream tasks like
humanization (which modifies framework residues).

\textbf{Why not learn the weights?}
Making weights trainable couples the masker to the model and
complicates reproducibility. Fixed weights keep the ablation clean.

\textbf{What if CDR annotation fails?}
The strategy silently falls back to uniform masking. This ensures
training never crashes due to missing metadata, and the model
still gets a gradient signal from that sample.

% ============================================================================
\section{Strategy 3: Span Masking}
\label{sec:span}
% ============================================================================

\subsection{The Problem}
When you mask random individual positions, predicting $x_i$ is
often trivially easy because $x_{i-1}$ and $x_{i+1}$ are both
visible. Protein sequences have strong local patterns (e.g.\
helices repeat every 3.6 residues), so single-position masking
gives away too much.

\subsection{The Solution}
Mask \textbf{contiguous spans} of residues instead of scattered
single positions. This forces the model to use longer-range
context. It also mimics the real task of designing a complete
loop or motif.

\subsection{How Span Lengths Are Chosen}
We draw span lengths from a \textbf{geometric distribution} with
parameter $p = 0.2$, capped at 10:

\begin{itemize}
\item Average span length: $1/0.2 = 5$ residues
\item Most spans are short (1--4 residues), but some reach 8--10
\item The cap at 10 prevents absurdly long spans that would swallow
      entire CDRs
\end{itemize}

Probability of each span length:
\begin{center}
\begin{tabular}{lccccccccc}
\toprule
Length & 1 & 2 & 3 & 4 & 5 & 6 & 7 & 8 & 9--10 \\
\midrule
Prob. & 20\% & 16\% & 13\% & 10\% & 8\% & 7\% & 5\% & 4\% & 6\% \\
\bottomrule
\end{tabular}
\end{center}

\subsection{Algorithm}
\begin{algorithm}[H]
\caption{Span Masking}
\begin{algorithmic}[1]
\State Compute budget: $T = \text{round}(0.15 \times \text{num maskable positions})$
\State $\text{masked\_count} = 0$
\While{$\text{masked\_count} < T$ and attempts $< 10T$}
    \State Pick a random unmasked, non-special position as the span start
    \State Draw span length from Geometric($0.2$), cap at 10
    \State Extend the span rightward from the start
    \State \textbf{Stop extending} if we hit a special token or end of sequence
    \State Mark newly-covered positions as masked
\EndWhile
\end{algorithmic}
\end{algorithm}

\subsection{Design Decisions}

\textbf{Why geometric distribution and not fixed-length spans?}
A mix of short and long spans gives the model diverse challenges.
Short spans test local prediction; long spans test long-range
reasoning. A fixed length would train only one difficulty level.

\textbf{Why stop at special tokens?}
In paired VH+VL sequences, a separator token marks the chain
boundary. Letting a span cross the separator would jumble two
chains together, which is biologically meaningless.

\textbf{Why cap at 10?}
A typical CDR3 is 10--15 residues. Capping at 10 means a span
can mask most of a CDR3 but not all of it plus its flanking
framework, preserving some local context.

\textbf{Why the $10T$ attempt limit?}
Safety valve. On very short sequences where almost everything is
already masked, without this cap the loop could run forever trying
to find an unmasked starting position.

\textbf{Why not sample spans centered on CDR regions?}
That would conflate span masking with CDR masking. We keep them
separate so the hybrid strategy can mix them independently.

% ============================================================================
\section{Strategy 4: Structure-Aware Masking}
\label{sec:structure}
% ============================================================================

\subsection{The Problem}
Sequence distance $\neq$ spatial distance. Two residues 30 positions
apart in sequence can be physical neighbors in the 3D fold (e.g.\
$\beta$-strands that pack against each other). Uniform and span
masking know nothing about this.

\subsection{The Solution: Seed-and-Grow on a 3D Neighborhood Graph}
We precompute a ``who-is-near-whom in 3D'' graph for each sequence,
then mask clusters of spatially-close residues.

\subsubsection{Step 1: Build the Neighborhood Graph (Offline)}
For each training sequence:
\begin{enumerate}
\item Run \textbf{ESM2} (a large protein language model) to get
      predicted contact probabilities --- an $L \times L$ matrix
      where entry $(i,j)$ is the probability that residues $i$ and
      $j$ are within 8\AA{} in the 3D structure.
\item For each residue $i$, keep its \textbf{top 32 neighbors}
      (highest contact probability). Store as a $L \times 32$ index
      table.
\end{enumerate}

This is done once before training; the training code never touches
ESM2.

\subsubsection{Step 2: Mask Spatial Patches (Per Sample)}
\begin{algorithm}[H]
\caption{Structure Masking: Seed-and-Grow}
\begin{algorithmic}[1]
\State Compute budget $T = \text{round}(0.15 \times \text{num maskable positions})$
\State $\text{masked\_count} = 0$
\While{$\text{masked\_count} < T$}
    \State Pick a random unmasked position as the \textbf{seed}
    \State Mask the seed
    \For{each of the seed's 32 nearest 3D neighbors (closest first)}
        \If{budget full} \textbf{break} \EndIf
        \If{neighbor not already masked}
            \State Mask it; increment count
        \EndIf
    \EndFor
\EndWhile
\end{algorithmic}
\end{algorithm}

This produces \textbf{3D patches}: clusters of residues that are
close in space but may be far apart in sequence.

\subsection{Design Decisions}

\textbf{Why 32 neighbors?}
ProteinMPNN (the best inverse-folding model) found that structural
context saturates around 32--48 nearest C$\alpha$ neighbors. Fewer
misses important contacts; more adds noise from distant residues.

\textbf{Why ESM2 contacts instead of real 3D structures?}
Real structures (AlphaFold, IgFold) take minutes per sequence and
don't scale to 500k training sequences. ESM2 contacts take seconds
per batch and are good enough --- they correlate well with true
contact maps for antibodies.

\textbf{Why seed-and-grow instead of one big cluster?}
One cluster would mask a single surface patch and leave the rest
of the protein untouched. Multiple seeds spread the mask across
different parts of the structure, giving broader coverage over
training.

\textbf{Why visit neighbors closest-first?}
The most tightly-coupled residues (highest contact probability)
should be masked first --- this creates the hardest prediction
challenge by removing the most informative context.

\textbf{What if we don't have structure data for a sequence?}
Falls back to uniform masking. The strategy logs how often this
happens so we can catch systematic data issues.

\textbf{Why is there also a C$\alpha$-coordinate code path?}
Legacy support. An earlier version used IgFold-predicted
C$\alpha$ coordinates and computed distances with
\texttt{torch.cdist}. The ESM2-contact path is faster and is
the main pipeline. If both exist, the kNN path wins.

% ============================================================================
\section{Strategy 5: Interface / Paratope Masking}
\label{sec:interface}
% ============================================================================

\subsection{The Problem}
The \textbf{paratope} is the set of residues that physically
contact the antigen --- these are the residues that \emph{define}
what the antibody does. They are only about 15--25\% of the
sequence. Uniform masking under-serves them.

\subsection{The Solution}
Give paratope residues higher masking weight, similar to CDR
masking but using binding-site data instead of region labels.

\subsection{Where Do Paratope Labels Come From?}
We cannot know the true paratope without a crystal structure of
the antibody-antigen complex (which we don't have for 500k
sequences). Instead, we estimate it:

\begin{enumerate}
\item Take the \textbf{TDC SAbDab\_Liberis} dataset: $\sim$1,000
      antibody sequences with experimentally-determined paratope
      positions from crystal structures.
\item Annotate each sequence with IMGT CDR regions using ANARCI.
\item For every (amino acid, region) pair, count how often that
      combination is a paratope residue. This gives us a lookup
      table:
      \[
      P(\text{paratope} \mid \text{amino acid}, \text{region})
      \]
      For example: Tyrosine in CDR3 $\approx 45\%$ paratope.
      Alanine in framework $\approx 3\%$ paratope.
\item For each training sequence, look up each residue's
      paratope probability. These are \textbf{soft labels}
      (floats between 0 and 1), not hard yes/no calls.
\end{enumerate}

\subsection{The Weighting Formula}
For each position $i$ with paratope probability $\pi_i$:
\[
\text{weight}_i = 1 + 5 \times \pi_i
\]
\begin{itemize}
\item Non-paratope ($\pi = 0$): weight = 1
\item Definite paratope ($\pi = 1$): weight = 6
\item Somewhere in between: linearly interpolated
\end{itemize}

This 6:1 ratio concentrates about 50--60\% of the mask budget
on paratope positions.

\subsection{Algorithm}
\begin{algorithm}[H]
\caption{Interface Masking}
\begin{algorithmic}[1]
\Require Soft paratope labels $\pi_i \in [0,1]$ for each position
\If{no paratope labels available}
    \State fall back to uniform masking
\EndIf
\State $w_i = 1 + 5 \times \pi_i$ for each position
\State Zero out weights at special tokens
\State Compute average weight; normalize probabilities (Section~\ref{sec:weighted-flip})
\State Flip weighted coins
\end{algorithmic}
\end{algorithm}

\subsection{Design Decisions}

\textbf{Why soft labels instead of hard 0/1?}
Hard paratope labels require a crystal structure of every
antibody-antigen complex --- impossible at scale. Soft labels
approximate paratope tendency from sequence alone and can be
computed for every training sequence.

\textbf{Why a (AA, region) table and not just the region?}
Different amino acids at the same position have very different
paratope probabilities. Tyrosine and tryptophan are common in
paratopes (aromatic stacking); glycine and alanine are rare.
The per-amino-acid resolution captures this.

\textbf{Why the minimum-count threshold of 5?}
Rare (AA, region) pairs (e.g.\ Cysteine in CDR3) would have
noisy estimates. If we have fewer than 5 observations, we fall
back to the region-average rate. This is a standard smoothing
trick.

\textbf{Why 6:1 ratio?}
Same reasoning as CDR3's weight: strong enough to shift
$\sim$50\% of the budget to paratope, but not so extreme that
non-paratope positions are starved.

% ============================================================================
\section{Strategy 6: Germline / Mutation-Aware Masking}
\label{sec:germline}
% ============================================================================

\subsection{The Problem: 85\% of Masking Is Wasted}

This is arguably the most important strategy. Here's the issue:

An antibody starts as a \textbf{germline} sequence (encoded in
your DNA), then undergoes \textbf{somatic hypermutation (SHM)}
in the immune response to improve binding. In a typical VH
sequence:
\begin{itemize}
\item $\sim$85\% of residues \textbf{match germline} exactly
      (never mutated)
\item $\sim$15\% have been \textbf{mutated by SHM} (the
      interesting ones)
\end{itemize}

With uniform masking, 85\% of the gradient signal goes to
teaching the model to predict \emph{germline} residues --- which
are boring, because the V gene identity largely determines them.
The functionally interesting SHM mutations get only 15\% of the
budget. This is the ``germline bias'' problem identified in the
AbLang-2 paper.

\subsection{The Solution}
Give mutated positions higher masking weight so they get more
training signal.

\subsection{How We Compute Germline Labels}

For each training sequence, we need to know which positions match
germline and which were mutated. We compute this \textbf{without
any external germline database}:

\begin{enumerate}
\item \textbf{Build consensus germline from data.} Group all
      training sequences by their V-gene call (e.g.\ IGHV1-2*01).
      For each V gene with $\geq$20 sequences, compute the
      \textbf{per-position consensus} (most common amino acid at
      each position). Because SHM is sparse ($\sim$15\% rate),
      the consensus is virtually identical to the true germline.
\item \textbf{Do the same for J genes}, using the FR4 segment
      after CDR3.
\item \textbf{Compare each sequence to its consensus:}
      \begin{itemize}
          \item Position matches consensus $\to$ label = 0
                (germline)
          \item Position differs from consensus $\to$ label = 1
                (mutated)
          \item CDR3 junction positions $\to$ label = 0.5
                (ambiguous)
      \end{itemize}
\end{enumerate}

\subsection{Why CDR3 Gets Label 0.5}
CDR3 is created by V(D)J recombination with random nucleotide
insertions --- it doesn't come from a single germline template.
There is no meaningful ``germline'' to compare against:
\begin{itemize}
\item Label = 1 would over-focus on CDR3 (duplicating CDR masking)
\item Label = 0 would ignore CDR3 entirely
\item Label = 0.5 gives it moderate weight --- a compromise
\end{itemize}

\subsection{The Weighting Formula}
For each position $i$ with germline label $g_i \in \{0, 0.5, 1\}$:
\[
\text{weight}_i = 1 + 5 \times g_i
\]
\begin{center}
\begin{tabular}{lcc}
\toprule
Position type & Label & Weight \\
\midrule
Germline match & 0 & 1 \\
CDR3 junction & 0.5 & 3.5 \\
SHM mutation & 1 & 6 \\
\bottomrule
\end{tabular}
\end{center}

\subsection{Quantitative Impact}
With 85\% germline (weight 1) and 15\% mutated (weight 6):
\[
\text{average weight} = 0.85 \times 1 + 0.15 \times 6 = 1.75
\]
\[
\text{fraction of budget on mutations} = \frac{0.15 \times 6}{1.75} \approx 51\%
\]
This is a \textbf{3.4$\times$ boost} over the 15\% that uniform
masking would give mutations.

\subsection{Algorithm}
\begin{algorithm}[H]
\caption{Germline Masking}
\begin{algorithmic}[1]
\Require Germline labels $g_i \in [0, 1]$ for each position
\If{no germline labels available}
    \State fall back to uniform masking
\EndIf
\State $w_i = 1 + 5 \times g_i$
\State Zero out at special tokens; normalize; flip weighted coins
\end{algorithmic}
\end{algorithm}

\subsection{Design Decisions}

\textbf{Why build germline from data instead of downloading the
real germline sequences?}
Real germline databases (IMGT) are at the nucleotide level and
would require translation, indel handling, and species matching.
The data-derived consensus is automatically compatible with our
tokenization and, with $\geq$20 sequences per gene, is
statistically indistinguishable from the true germline.

\textbf{Why require $\geq$20 sequences per V gene?}
With fewer sequences, a single bad record could flip the
consensus. Sequences from rare V genes get label 0.5 everywhere
and effectively receive uniform masking.

\textbf{Why not a more extreme ratio like 100:1?}
At 100:1, virtually all masking would land on the $\sim$18
mutated positions in a 120-residue sequence. The model would
\emph{only} learn to predict mutations and lose all ability to
reconstruct framework --- bad for tasks like humanization or
stability prediction.

% ============================================================================
\section{Strategy 7: Multispecific / Paired-Chain Masking}
\label{sec:multispecific}
% ============================================================================

\subsection{The Problem}
The strategies above all work on single chains. But real antibodies
are \textbf{pairs}: a heavy chain (VH) and a light chain (VL)
that work together. For paired and multispecific antibodies, three
new challenges arise:
\begin{enumerate}
\item Which VL pairs with which VH matters for binding.
\item In bispecific antibodies, the light chain is often
      \textbf{shared} between two different heavy chains.
\item A small set of \textbf{interface residues} at the VH--VL
      packing surface determines pairing geometry.
\end{enumerate}

\subsection{The Solution: Three Policies, Randomly Mixed}
Each training step, we randomly pick one of three policies:

\subsubsection{Policy A: Module-Isolated Paratope Masking}

\textbf{Idea:} Pick one module (e.g.\ the VH of arm 1), mask
mostly its paratope residues, but let $\sim$10\% of the mask
budget ``leak'' to the other module.

\begin{itemize}
\item Within the target module: weight positions by paratope
      probability (same as interface masking, weight $= 1 + 5\pi$)
\item Other module: weight = $0.1$ (a small leak)
\end{itemize}

\textbf{Why the 10\% leak?}
Without it, the non-target module \emph{never} gets masked.
Over many steps, this would cause its representation to collapse
(no gradient signal $\to$ no learning). The leak keeps it alive.

\subsubsection{Policy B: Shared Light Chain Boost}

\textbf{Idea:} Mask the light chain 3$\times$ more heavily than
the heavy chain. This trains the model to predict VL conditioned
on VH --- exactly the task for shared-light-chain bispecific design.

\begin{center}
\begin{tabular}{lc}
\toprule
Chain & Weight \\
\midrule
Heavy (VH) & 1 \\
Light (VL) & 3 \\
\bottomrule
\end{tabular}
\end{center}

With roughly equal chain lengths, this puts $\sim$75\% of the
mask budget on the light chain.

\subsubsection{Policy C: VH--VL Interface Masking}

\textbf{Idea:} Focus masking on the $\sim$8 canonical interface
positions per chain (from Chothia \& Lesk, 1985) that form the
physical VH--VL packing surface. Masking these forces the model
to use the opposite chain as context --- the cleanest signal for
learning chain pairing.

Interface labels:
\begin{itemize}
\item Canonical interface positions: label $= 1.0$
\item Immediately adjacent positions ($\pm 1$): label $= 0.3$
\item Everything else: label $= 0$
\end{itemize}

Weights: $w_i = 1 + 5 \times \text{interface\_label}_i$, same
formula as the other weighted strategies.

\textbf{Why the 0.3 neighbor label?} The canonical positions alone
are only $\sim$8 per chain --- too sparse to generate meaningful
mask signal. Including neighbors at reduced weight smooths the
signal and reflects the physical reality that packing involves a
continuous contact surface, not isolated atoms.

\subsection{Policy Mixing}
Each training step, sample one policy from
$(\frac{1}{3}, \frac{1}{3}, \frac{1}{3})$. Default is equal
weighting (tunable via config).

\subsection{Design Decisions}

\textbf{Why three separate policies instead of one combined
weighting?}
Combining all three signals into one weight vector would blur
them together --- the gradient signal becomes an uninterpretable
average. By committing to one policy per sample, each gradient
step has a clear learning objective. Over a batch of 32 samples,
roughly 10--11 use each policy.

\textbf{Why does Policy B boost light chain and not heavy?}
In bispecific antibodies, the light chain is the \emph{shared}
component. Training the model to reconstruct VL given VH is
the direct analogue of the downstream design task. The reverse
direction is already covered by Policy A.

% ============================================================================
\section{Strategy 8: Hybrid (The Meta-Strategy)}
\label{sec:hybrid}
% ============================================================================

\subsection{The Problem}
Each strategy above captures one aspect of antibody biology.
We want a model that understands \emph{all of them}. Simply
picking one strategy for the whole training run would miss the
others.

\subsection{The Solution: Mix Strategies Per Sample}
The hybrid strategy holds a collection of sub-strategies and,
for each training sample, randomly picks one according to a
mixture distribution:

\begin{enumerate}
\item Check which sub-strategies have the required metadata for
      this particular sample (e.g.\ skip germline masking if this
      sample has no germline labels).
\item Renormalize the mixture weights over only the available
      strategies.
\item Sample one strategy and delegate to it.
\end{enumerate}

\subsection{Curriculum Scheduling}
The mixture weights can change over training via a
\textbf{piecewise-linear schedule}:

\begin{center}
\small
\begin{tabular}{@{}r r r r r r r l@{}}
\toprule
Step & uniform & cdr & span & struct. & interf. & germ. & Rationale \\
\midrule
0 & 0.30 & 0.15 & 0.30 & 0.10 & 0.10 & 0.05 & Learn basic protein syntax \\
6k & 0.10 & 0.30 & 0.10 & 0.15 & 0.25 & 0.10 & Specialize on binding sites \\
19k & 0.10 & 0.20 & 0.10 & 0.15 & 0.20 & 0.25 & Learn mutation patterns \\
40k+ & 0.10 & 0.20 & 0.10 & 0.15 & 0.20 & 0.25 & Freeze (stable training) \\
\bottomrule
\end{tabular}
\end{center}

Between breakpoints, weights are \textbf{linearly interpolated}
(smooth transition, no sudden jumps that could destabilize
training).

\textbf{Stage logic:}
\begin{itemize}
\item \textbf{Stage 1} (steps 0--6k): Heavy on uniform + span.
      The model is still warming up (this coincides with the
      learning rate warmup). Easy masking for a fragile model.
\item \textbf{Stage 2} (6k--19k): Shift to CDR + interface.
      The model is stable enough for harder, biologically-focused
      masking.
\item \textbf{Stage 3} (19k--40k): Increase germline weight.
      Now the model is sophisticated enough to benefit from the
      subtle SHM signal.
\item \textbf{Stage 4} (40k+): Freeze. Stable training for the
      remaining $\sim$68\% of training at a balanced mixture.
\end{itemize}

\subsection{Important Implementation Detail: Shared Memory}

When training uses multiple data-loader workers (which we do for
speed), each worker gets a \emph{copy} of the strategy object.
If we naively update the mixture weights in the main process, the
worker copies never see the change --- the curriculum silently
does nothing.

\textbf{Fix:} We store the weight tensor in shared memory
(\texttt{share\_memory\_()}) and update it in-place
(\texttt{copy\_()}). All worker processes reference the same
underlying memory, so they all see curriculum updates immediately.

This is a subtle but critical correctness issue --- without it,
the curriculum appears to work (the main process logs show updated
weights) but has zero effect on actual masking.

\subsection{Algorithm}
\begin{algorithm}[H]
\caption{Hybrid Masking}
\begin{algorithmic}[1]
\State \textbf{Setup:} Initialize all sub-strategies; store weights in shared memory
\State \textbf{Each training step:} Callback interpolates curriculum $\to$ updates shared weights
\State \textbf{Each sample:}
\State \hspace{1em} Find which sub-strategies have metadata for this sample
\State \hspace{1em} Renormalize mixture weights over available strategies
\State \hspace{1em} Sample one strategy from the renormalized distribution
\State \hspace{1em} Delegate to that strategy's \texttt{select\_mask\_positions()}
\end{algorithmic}
\end{algorithm}

\subsection{Design Decisions}

\textbf{Why sample one strategy per sample instead of per batch?}
Per-batch means all 32 samples use the same strategy ---
gradients within a batch are correlated. Per-sample decorrelates
them, giving the optimizer a cleaner gradient estimate.

\textbf{Why not apply all strategies simultaneously (union of masks)?}
That would mask far more than 15\% of positions and confound the
budget. The mixture sampler is an unbiased estimator of the
combined loss at 1$\times$ cost.

\textbf{Why include uniform in the mixture at all?}
It guarantees a fallback: even if all metadata is missing, the
uniform sub-strategy is always available.

\textbf{Why renormalize over available strategies?}
Without renormalization, a sample missing one strategy's metadata
would have ``dead'' probability mass, biasing toward fewer total
masks. Renormalization ensures every sample fully uses the budget.

\textbf{Why linear interpolation between curriculum breakpoints?}
Step functions cause sudden distribution shifts that can spike the
training loss. Linear interpolation is smooth and easy to reason
about.

\textbf{Why freeze after step 40k?}
Continuing to shift weights late in training introduces
distributional drift when we want the model to converge.

% ============================================================================
\section{How Everything Fits Together}
% ============================================================================

\subsection{Architecture Overview}
\begin{enumerate}
\item \textbf{Config} (YAML): specifies which strategy to use
      and all hyperparameters.
\item \textbf{Registry}: strategies register themselves by name;
      training code looks them up by string.
      Adding a new strategy $=$ one Python file, zero changes to
      training code.
\item \textbf{Collator}: pads sequences, builds per-sample
      metadata, calls the strategy. Strategy-agnostic --- swap
      the strategy and nothing else changes.
\item \textbf{Two-stage design}: every strategy only implements
      ``pick positions.'' The 80/10/10 corruption is shared.
\end{enumerate}

\subsection{Summary Table}

\begin{center}
\renewcommand{\arraystretch}{1.3}
\small
\newcolumntype{L}{>{\raggedright\arraybackslash}X}
\begin{tabularx}{\textwidth}{@{}l L L@{}}
\toprule
\textbf{Strategy} & \textbf{What it focuses on} & \textbf{Key idea} \\
\midrule
\strategy{uniform} & Nothing (equal everywhere) & Baseline control \\
\strategy{cdr} & CDR loops (especially CDR3) & Weight CDRs 3--6$\times$ higher \\
\strategy{span} & Contiguous chunks & Geometric-length spans \\
\strategy{structure} & 3D spatial neighborhoods & Seed-and-grow on kNN graph \\
\strategy{interface} & Paratope (binding) residues & Soft labels from TDC data \\
\strategy{germline} & SHM mutations & 3.4$\times$ boost to mutated positions \\
\strategy{multispecific} & VH--VL pairing & Three-policy mixture (A/B/C) \\
\strategy{hybrid} & All of the above & Per-sample mixture with curriculum \\
\bottomrule
\end{tabularx}
\end{center}

% ============================================================================
\section{Anticipated Questions}
% ============================================================================

\textbf{Q: Why is the mask rate always 15\%?}

So comparisons are fair. If two strategies use different rates,
any difference in performance could be from the rate, not the
strategy. The mask rate is a separate knob to tune.

\textbf{Q: What happens when metadata is missing?}

Every strategy falls back to uniform masking. Training never
crashes. Fallback rates are logged so we catch systematic issues.

\textbf{Q: Why not one ``super-weight'' that combines CDR +
germline + paratope + structure?}

Multiplying weights ($w_{\text{cdr}} \times w_{\text{germ}}$)
produces extreme spikes; adding them is dominated by whoever has
the largest magnitude. The hybrid sampler avoids this by
committing to one clean signal per sample.

\textbf{Q: Why not learn the mixture weights?}

Learned weights create a feedback loop: if the model is bad at
paratope prediction, it would down-weight paratope masking,
further starving paratope learning. Fixed/scheduled weights
maintain a stable learning target.

\textbf{Q: Why does the curriculum start easy?}

The learning rate is still warming up in the first 6k steps.
Hitting a fragile, randomly-initialized model with hard masking
signals (germline, structure) before it has learned basic protein
syntax is counterproductive.

\textbf{Q: What if a sequence has no CDR annotations at all?}

The CDR weights all become 1 (framework), so the formula
$0.15 \times 1/1 = 0.15$ gives exact uniform masking. The
strategy degrades to the baseline automatically --- this is a
built-in property of the weighted coin-flip formula.

\textbf{Q: Why does structure masking not also use sequence
distance?}

That is exactly what span masking does. Structure masking is the
\emph{complement} --- coupling residues that are close in 3D but
far in sequence. Combining them would collapse toward uniform.

\textbf{Q: Is the soft paratope table biased because TDC is small?}

Yes, acknowledged limitation. The table is smoothed (region-level
fallback for rare cells), and the labels are used as
\emph{weights}, not ground truth. Biased weights still
concentrate gradient on the right neighborhood of positions, even
if the exact probabilities are off.

\textbf{Q: How does the shared-memory fix work exactly?}

\texttt{torch.Tensor.share\_memory\_()} moves the tensor's
storage to \texttt{/dev/shm} (POSIX shared memory). When the
main process calls \texttt{tensor.copy\_(new\_values)}, it writes
to that shared memory region. Forked workers hold references to
the same region, so they read the updated values without any
explicit synchronization.

\end{document}
